\documentclass[10pt,a4paper]{amsart}
\usepackage[margin=19mm]{geometry}
\usepackage{amsmath,amssymb,mathtools}
\usepackage[scr=rsfs,cal=euler]{mathalfa}
\usepackage{xfrac}
\usepackage[unicode,hidelinks,bookmarks=false]{hyperref}
\newcommand{\ie}{i.e.\ }
\newcommand{\cf}{cf.\ }
\newcommand{\resp}{respectively\ }
\newcommand{\Subsection}{\subsection}
\providecommand{\nobreakdash}{\nobreak\hbox{-}}
\setlength{\emergencystretch}{2em}
\multlinegap0pt
\usepackage{mathtools}
\usepackage{tikz}
\usepackage{algorithm}\usepackage{algpseudocode}
\usepackage[normalem]{ulem}
\usepackage{cancel}
\usepackage{enumerate}
\usepackage{amssymb}
\newtheorem{proposition}{Proposition}
\newtheorem{lemma}{Lemma}
\newcommand{\cS}{\mathcal{S}}
\newcommand{\one}{\mathbf 1}
\title{An Exact Pointwise Characterization for Total Variation Denoising in Quantile Regression}
\author{Deep Ghoshal, Sabyasachi Chatterjee}
\date{Source: arXiv:2605.01237v1 (CC BY 4.0)}
\begin{document}
\maketitle

\noindent\emph{Source and licence.} An Exact Pointwise Characterization for Total Variation Denoising in Quantile Regression. Version arXiv:2605.01237v1 (CC BY 4.0), distributed by arXiv under the Creative Commons Attribution 4.0 International licence, http://creativecommons.org/licenses/by/4.0/, as recorded in the arXiv licence metadata for this exact version and shown on the abstract page https://arxiv.org/abs/2605.01237v1. Full licence notice: \texttt{licenses/arxiv-2605.01237-CC-BY-4.0.txt}.\par\smallskip

\noindent\textit{Editorial scope.} Counted unit: the article's proposition lem:perturb-upper (source0.tex lines 738--753). The article's complete original proof of it is delivered with it (source0.tex lines 754--863, one \texttt{proof} environment of 110 lines). No mathematics is written, repaired or re-derived: the statement and the proof are the article's own lines, in the article's own notation, and no retained line is substituted or re-flowed.  No further source-local context is needed: every cross-reference of every retained line resolves inside the two delivered spans.  Nothing was omitted from any retained span, so the package carries no omission record.  No retained line contains a citation, so the wrapper carries no bibliography and no bibliography entry.  No retained line contains a figure environment, an image-inclusion command or drawing code, so no image is delivered and the package declares no asset dependency.  Editorial formatting only, in addition: an AMS \texttt{amsart} preamble that restates the article's own declarations and macros used by the retained lines, namely the environments \texttt{proposition}, \texttt{lemma} on separate counters, together with the standard packages those lines need.  The retained environments print in this document as Proposition 1, Lemma 1 in order of appearance, whereas the article prints the same items with the numbering of its own full document.  The complete native source of the article as distributed in the arXiv e-print for this exact version is bundled unchanged as \texttt{sources/199621/source0.tex}.
\begin{proposition}[One order statistic upgrade]\label{lem:perturb-upper}
 Fix any location $i \in [n]$ and define
    \[\hat\theta_i^{\max}:=\max\{\theta_i:\theta\in\mathcal \cS\}.\]
    Moreover, let $\hat{\theta}^{\max}$ be an element in $\cS$ corresponding to $\hat{\theta}^{\max}_i$.
Let $\hat{J}=[a:b]$ be the largest interval containing $i$ such that
\[
\hat\theta^{\max}_u \le \hat\theta^{\max}_i
\qquad \forall u\in \hat J.
\]
Fix any $I=[c:d]\subseteq \hat J$ with $i\in I$ and suppose
$u_{I,\hat J}\in\mathbb Z$ and $0\le u_{I,\hat J}<|I|$.
Then,
\[
\hat\theta_i^{\max}\ge y_{I,(u_{I,\hat J}+1)}.
\]
\end{proposition}

\begin{proof}[Proof of Proposition~\ref{lem:perturb-upper}]
Let $K:=u_{I,\hat J}$ and set
\[
q:=y_{I,(K+1)}.
\]
Suppose for contradiction that $\hat\theta_i^{\max}<q$.
Define
\begin{equation}\label{defn:theta_t}
    \theta(t):=\hat\theta^{\max}+t\,\one_I,\qquad t\ge 0,
\end{equation}

and $\phi(t):=F(\theta(t))$. We will argue that $\phi$ is affine in some $[0,\eta_1]$ for some small enough $\eta_1>0$ and $\phi'_{+}(0)\leq 0$. To that end, for $j\in I$ write
\[
u_j(t):=y_j-\theta_j(t)=y_j-\hat\theta^{\max}_j-t.
\]
If $y_j>\hat\theta^{\max}_j$, then $u_j(0)>0$ and the sign remains positive for
$t<y_j-\hat\theta^{\max}_j$.
If $y_j=\hat\theta^{\max}_j$, then $u_j(t)=-t<0$ for all $t>0$.
If $y_j<\hat\theta^{\max}_j$, then $u_j(t)<0$ for all $t\ge 0$.
Hence, letting
\[
\eta:=\min_{j\in I:\,y_j>\hat\theta^{\max}_j}
(y_j-\hat\theta^{\max}_j)>0,
\]
(with $\eta=+\infty$ if the index set is empty), we conclude that $u_j(t)$ does not change sign for $t\in [0,\eta]$ for any $j \in I.$ Thus, from the definition of $\rho_{\tau}$, it follows that each $\rho_\tau(u_j(t))$ is affine on $[0,\eta]$,
and so is their sum. Thus,
\[
\frac{d}{dt}\rho_\tau(u_j(t))\Big|_{0^+}
=
\begin{cases}
-\tau, & \hat\theta^{\max}_j<y_j,\\
1-\tau, & \hat\theta^{\max}_j\ge y_j.
\end{cases}
\]
Summing gives
\[
\frac{d}{dt}\sum_{j=c}^d \rho_\tau(y_j-\theta_j(t))\Big|_{0^+}
=
-\tau|I|+\#\{j\in I:\hat\theta^{\max}_j\ge y_j\}.
\]

Now, we turn our attention to the penalty term. Only the two boundary edges across $\partial \hat{J}$ depend on $t$ in the TV term. In the following lemma, we establish that $TV(\theta(t))$ is affine in a neighborhood around $0$.
\begin{lemma}\label{lemma:TV_is_affine}
    Let $\theta(t)$ be as defined in~\eqref{defn:theta_t}. Then, for small enough $t>0$, it holds that
    \begin{equation}\label{eqn:penalty_term_is_affine}
    TV(\theta(t))-TV(\theta(0))=(\sigma_{c-1}\one\{c\neq 1\}-\sigma_d\one\{d\neq n\})t,
\end{equation}
where for $c>1$ and $d<n$,
\[
\sigma_{c-1}:=\operatorname{sign}_+\bigl(\hat\theta^{\max}_c-\hat\theta^{\max}_{c-1}\bigr),
\qquad
\sigma_d:=\operatorname{sign}_-\bigl(\hat\theta^{\max}_{d+1}-\hat\theta^{\max}_d\bigr),
\]
and $\operatorname{sign}_+,\;\operatorname{sign}_-$ are defined as
\[
\operatorname{sign}_+(x):=
\begin{cases}
+1,& x\ge 0,\\
-1,& x<0,
\end{cases}
\qquad
\operatorname{sign}_-(x):=
\begin{cases}
+1,& x>0,\\
-1,& x\le 0.
\end{cases}
\]
\end{lemma}
Hence, we conclude that
\[\frac{d}{dt}\Big(\lambda\sum_{k}|\theta_{k+1}(t)-\theta_k(t)|\Big)\Big|_{0^+}
=
\lambda(\sigma_{c-1}\one\{c\neq 1\}-\sigma_d\one\{d\neq n\}).\]
Thus, we conclude that there exists $0<\eta_1\leq \eta$ such that $\phi$ is affine in $[0,\eta_1]$ and 
\begin{equation}\label{eqn:phi_deriv}
    \phi'_+(0)=-\tau|I|+\#\{j\in I:\hat\theta^{\max}_j\ge y_j\}+\lambda(\sigma_{c-1}\one\{c\neq 1\}-\sigma_d\one\{d\neq n\}).
\end{equation}
Now, since $I\subseteq J$ and $\hat\theta^{\max}_j\le\hat\theta^{\max}_i$ for all $j\in I$,
\[
\#\{j\in I:\hat\theta^{\max}_j\ge y_j\}
\le
\#\{j\in I:y_j\le \hat\theta_i^{\max}\}.
\]
However, because we assumed $\hat\theta_i^{\max}<q=y_{I,(K+1)}$,
\[
\#\{j\in I:y_j\le \hat\theta_i^{\max}\}\le \#\{j\in I:y_j< y_{I,(K+1)}\}= K.
\]
Thus,
\begin{equation}\label{eq:data-bound-full}
-\tau|I|
+\#\{j\in I:\hat\theta^{\max}_j\ge y_j\}
\leq
-\tau|I|+K
=
-2\lambda C_{I,\hat{J}},
\end{equation}
where the last equality follows from the definition of $u_{I,\hat{J}}=K.$
Now, in the following lemma, we establish an upper bound on the third term of~\eqref{eqn:phi_deriv}.
\begin{lemma}\label{lemma:tv-bound-full-lemma}
    \begin{equation}\label{eq:tv-bound-full-lemma}
\lambda(\sigma_{c-1}\one\{c\neq 1\}-\sigma_d\one\{d\neq n\})\le 2\lambda C_{I,J},
\end{equation}
where the associated quantities are defined in Lemma~\ref{lemma:TV_is_affine}.
\end{lemma}
Plugging in \eqref{eq:data-bound-full} and \eqref{eq:tv-bound-full-lemma} in \eqref{eqn:phi_deriv}, we conclude that $\phi'_+(0)\le 0$. However, since $\hat\theta^{\max}$ minimizes $F$, $\phi'_+(0)\ge 0$. Thus, it must holds that $\phi'_+(0)=0$.
However, we have already argued that $\phi$ is affine on $[0,\eta_1]$. Therefore, it must hold that $\phi(t)=\phi(0)$ for $t\in[0,\eta_1]$. Choose $\varepsilon\in(0,\eta_1]$ and define
$\tilde\theta=\hat\theta^{\max}+\varepsilon\one_I$.
Then $\tilde\theta\in\mathcal \cS$ and
$\tilde\theta_i>\hat\theta_i^{\max}$,
contradiction! Thus, the assumption that $\hat\theta_i^{\max}<q$ must be wrong and thus, the proof of Proposition~\ref{lem:perturb-upper} is complete.
\end{proof}

\end{document}
